\begin{table}[p]\centering\footnotesize
\setlength{\tabcolsep}{4pt}
\begin{tabular}{lrrrrr}\toprule
Feature & 0005 (A) & 0007 (B) & 0010 (B) & 0011 (DT-68) & 0012 (K) \\\midrule
\multicolumn{6}{l}{Selected amino-acid modules: completeness (\%)} \\
Cysteine from methionine & 33.3/33.3 & 33.3/33.3 & 33.3/33.3 & 33.3/33.3 & 16.7/16.7 \\
Methionine from aspartate & NR/14.3 & 14.3/14.3 & 14.3/14.3 & NR/NR & NR/NR \\
Histidine from PRPP & 5.6/11.1 & 5.6/5.6 & 5.6/5.6 & NR/8.3 & NR/NR \\
Arginine from ornithine & 33.3/33.3 & NR/NR & NR/NR & 33.3/33.3 & 33.3/33.3 \\
\midrule\multicolumn{6}{l}{Selected defense genes: assigned copies} \\
Cas1 & 0/0 & 0/0 & 0/0 & 0/0 & 0/0 \\
Cas2 & 0/0 & 0/0 & 0/0 & 0/0 & 0/0 \\
Cas9/Csn1 & 0/0 & 0/0 & 0/0 & 0/0 & 0/0 \\
Csn2 & 0/0 & 0/0 & 0/0 & 0/0 & 0/0 \\
Type I restriction R & 0/0 & 0/0 & 1/1 & 0/0 & 0/0 \\
Type I modification M & 0/0 & 0/0 & 1/1 & 0/0 & 0/0 \\
Type I specificity S & 0/0 & 0/0 & 0/1 & 0/0 & 0/0 \\
DNA methyltransferase & 0/0 & 1/1 & 0/0 & 0/0 & 0/0 \\
Type III restriction & 0/0 & 1/1 & 0/0 & 0/0 & 0/0 \\
\bottomrule\end{tabular}
\caption{Selected nutritional and defense annotations in near-complete Metamycoplasmataceae genomes. Column identifiers omit the MAGSP prefix; A, B, DT-68 and K identify the four focal lineages. Cells give strict/relaxed assignments. NR means no row was returned in the stored module output and is not an imputed completeness estimate. Module identifiers are M00609, M00017, M00026 and M00844, respectively; these overlapping and precursor-dependent routes cannot be counted as experimentally established auxotrophies. Defense entries are selected orthologs, not counts of functional systems. Incomplete assemblies can lose genes. Machine-readable tables additionally retain the two presence-only genomes, all selected amino-acid module records, missing orthologs, and the screened chitin-degrading genes.}
\label{tab:nutrition-defense}\end{table}
